\documentclass[10pt,a4paper]{amsart}
\usepackage[margin=19mm]{geometry}
\usepackage{amsmath,amssymb,mathtools}
\usepackage[scr=rsfs,cal=euler]{mathalfa}
\usepackage{xfrac}
\usepackage[unicode,hidelinks,bookmarks=false]{hyperref}
\newcommand{\ie}{i.e.\ }
\newcommand{\cf}{cf.\ }
\newcommand{\resp}{respectively\ }
\newcommand{\Subsection}{\subsection}
\providecommand{\nobreakdash}{\nobreak\hbox{-}}
\setlength{\emergencystretch}{2em}
\multlinegap0pt
\usepackage{bm}
\usepackage{bbm}
\usepackage{dsfont}
\usepackage{xspace}
\usepackage{cancel}
\usepackage{mathtools}
\usepackage{amsthm}
\makeatletter
\@ifundefined{theorem}{\newtheorem{theorem}{Theorem}}{}
\@ifundefined{lemma}{\newtheorem{lemma}[theorem]{Lemma}}{}
\makeatother
\newcommand{\Fq}{\mathbb{F}_q} 
\title[Combinatorial Bounds for Codes over Metric Spaces: Ramsey-Si]{Combinatorial Bounds for Codes over Metric Spaces: Ramsey-Sidorenko Thresholds and Subgraph Counts}
\author{Lucas Waite (Kenyon College), Nuh Aydin (Kenyon College)}
\date{Source: arXiv:2607.27098v1 (CC BY 4.0)}
\begin{document}
\maketitle

\noindent\emph{Source and licence.} Combinatorial Bounds for Codes over Metric Spaces: Ramsey-Sidorenko Thresholds and Subgraph Counts. Version arXiv:2607.27098v1 (CC BY 4.0), distributed under the Creative Commons Attribution 4.0 International licence, as recorded in the arXiv licence metadata for this exact version and shown on the abstract page https://arxiv.org/abs/2607.27098v1. Full rights notice: licenses/arxiv-2607.27098-CC-BY-4.0.txt.\par\smallskip

\noindent\textit{Editorial scope.} Counted unit: the Lemma whose source label is \texttt{lem: entropy-inequality} in the article (source0.tex lines 428--437, source label \texttt{lem: entropy-inequality}), delivered together with the article's own complete original proof of it (source0.tex lines 439--640, one \texttt{proof} environment of 202 lines). No mathematics is written, repaired or re-derived: statement and proof are the article's own text line for line. Required context retained uncounted: none. Every definition, display and label of those lines is retained unchanged. Omitted from this package and not redistributed: 1--427, 438--438, 641--1103. Nothing inside a retained excerpt is omitted or altered: every retained body is delivered exactly as the article's lines are written, so no omission receipt of any kind accompanies this package. Citations: the retained excerpts carry no citation command at all, so no bibliography entry of the article is transcribed and no reference is redistributed. Reference adaptation: none was needed and none was made; every reference inside the delivered unit resolves inside it, so no unresolved reference or citation is produced. Printed slips kept literal: no printed word, symbol, hypothesis or number of the counted statement or of its proof was altered. Layout and class: the article's own class is \texttt{article} with packages including graphicx, amsfonts, amsmath, amssymb, amsthm, newtxtext, xcolor, soul, fullpage, comment, hyperref; the wrapper is the lane's pinned \texttt{amsart} editorial wrapper at 10pt on A4, which restates the theorem-like environments the retained text needs and adds the article's own macro definitions that the retained text needs (\texttt{\textbackslash Fq}), with extra wrapper packages (bm, bbm, dsfont, xspace, cancel, mathtools). The delivered numbering is the unit's own. Assets: no retained span contains a figure environment, a graphics-inclusion command, a PDF-page inclusion, a table or a TikZ picture; no image file is copied into this package, the package declares no asset dependency and no image file is redistributed with it. Licence: version arXiv:2607.27098v1 of this article is distributed under the Creative Commons Attribution 4.0 International licence, as recorded in the arXiv licence metadata for this exact version and shown on the abstract page https://arxiv.org/abs/2607.27098v1. A full rights notice is delivered at licenses/arxiv-2607.27098-CC-BY-4.0.txt. Characters: every retained excerpt is pure ASCII, so the delivered engine, XeLaTeX with its default Latin Modern text font, reports no missing character.
\begin{lemma}
\label{lem: entropy-inequality}
Let $q$ be a prime power and $0<\delta<1-\frac1q$. For a given graph $H$, fix any vertex $r\in V(H)$ and let $\Omega=\{x:V(H)\to\Fq:x_r=0\}$. Then

\begin{align*}
    \Phi_{H,q}(\delta)&:= \max\left\{\mathcal H_q(\mu):\Pr_\mu(x_u\not=x_v)\leq\delta\text{ for every }e=uv\in E(H)\right\}\\
    &\ge v-1-m+mh_q(\delta).
\end{align*}

\end{lemma}

\begin{proof}
Throughout this argument we let $\varphi(x)=x\log_qx$.

First, our definition of $\Phi_{H,q}(\delta)$ as the maximum rather than the supremum is well-defined: Since the set of distributions $\mu$ satisfying $\Pr_\mu(x_u\not=x_v)\leq\delta$ for every $e\in E(H)$ forms a nonempty closed and bounded subset of the finite-dimensional space $\mathbb R^{|\Omega|}$ and base $q$ entropy is continuous, we can take a probability distribution $\mu$ on $\Omega$ attaining $\Phi_{H,q}(\delta)$.

Now we consider $\Phi_{H,q}$ as an optimization problem. Consider $\mu$ a vector in $\mathbb R^{|\Omega|}$, and let $a_{e}:=(\mathbf1_{\{x_u\not=x_v\}})_{x\in\Omega}$ for any $e=uv\in E(H)$. Then we find $-\Phi_{H,q}$ as
\[
\begin{aligned}
\min\quad & -\mathcal{H}_q(\mu)\\
\text{subject to}\quad
&\mathbf1^\top\mu-1=0,\\
&-\mu_x\leq0\qquad x\in\Omega,\\
&a_{e}^\top\mu-\delta\leq0\qquad e\in E(H)
\end{aligned}
\]
Each constraint is affine, and $-\mathcal{H}_q(\mu)$ is convex and differentiable over $\mu$ with full support, so we  show the optimal $\mu$ has full support in order to show we can use the KKT theorem:

Define $\nu$ coordinate-wise with $\nu_x=(1-\epsilon)\mathbf1_{\{x=\mathbf0\}}+\epsilon V_x$ where $V$ is the uniform distribution on $\Omega$ viewed as a vector in $\mathbb R^{|\Omega|}$. Then
\[
\Pr_{\nu}(x_u\not=x_v)=\epsilon\bigl(1-\frac 1q\bigr)<\delta\text{ choosing }0<\epsilon<\delta\frac{q}{q-1},
\]
and $\nu$ has full support: for any $x\in\Omega$, $\nu_x\geq\epsilon V_x>0$. Now we mix $\mu$ and $\nu$: let $\sigma(t)=(1-t)\mu+t\nu$ for $0<t<1$.

We now compare the contributions of $x$ yielding $\mu_x=0$ and $\mu_x>0$ to the difference in entropy of $\mu$ and $\sigma(t)$:

Let $P=\{x:\mu_x>0\}$ and consider $x\in P$. By the mean-value theorem,
\[
\varphi(\mu_x)-\varphi(\sigma(t)_x)=\varphi'(\beta_x(t))(\mu_x-\sigma(t)_x)=\varphi'(\beta_x(t))t(\mu_x-\nu_x)
\] 
for some $\beta_x(t)$ between $\mu_x$ and $\sigma(t)_x$. If $\nu_x\not=\mu_x$, we  let $t\leq\frac{\mu_x}{2|\nu_x-\mu_x|}$ for all of the finitely many $x\in P$ so that $\sigma(t)_x\in[\frac{\mu_x}2,\frac{3\mu_x}2]$, implying that
\[
|\varphi'(\beta(t)_x)|\leq\sup_{s\in[{\mu_x/}2,\ 3\mu_x/2]}\left|\varphi'(s)\right|<\infty.
\]
For the $x\in P$ with $\mu_x=\nu_x$ this follows immediately.

As a result, $\varphi(\mu_x)-\varphi(\sigma(t)_x)\leq C_xt$ for 
\[
C_x=|\mu_x-\nu_x|\sup_{s\in[{\mu_x/}2,\ 3\mu_x/2]}\left|\varphi'(s)\right|
\]
a constant with respect to $t$. 

Let $C=\sum_{x\in P}C_x$. Then for sufficiently small $t$,
\[
\sum_{x\in P}(\varphi(\sigma(t)_x)-\varphi(\mu_x))\geq-Ct.
\]

Now, if $x_0$ satisfies $\mu_{x_0}=0$, then
\begin{align*}
    \mathcal H_q(\sigma(t))-\mathcal H_q(\mu)&\geq\sum_{x\in P}\left(\varphi(\sigma(t)_x)-\varphi(\mu_x)\right)+\varphi(\sigma(t)_{x_0})\\
    &\geq-Ct+t\nu_{x_0}\log_q\frac{1}{t\nu_{x_0}}\\
    &\geq t(\nu_{x_0}\log_q\frac{1}{t}-C)\\
    &>0
\end{align*}
choosing $t$ sufficiently small.
This would therefore imply $\sigma(t)$ has higher entropy than $\mu$, and as $\sigma(t)$ is feasible, as it is a convex combination of feasible distributions, and $\mu$ is optimal, this cannot be the case. Thus, $\mu$ must have full support. We thus obtain the equivalent simplified optimization problem:

\[
\begin{aligned}
\min_{(0,\infty)^{|\Omega|}}\quad & -\mathcal{H}_q(\mu)\\
\text{subject to}\quad
&\mathbf1^\top\mu-1=0,\\
&a_{e}^\top\mu-\delta\leq0\qquad e\in E(H)
\end{aligned}
\]
Since the original optimizer \(\mu\) has full support, it is also a
local minimizer of this restricted problem. We also have shown $-\mathcal H_q$ is differentiable on an open neighborhood of $\mu$ contained in $(0,\infty)^{|\Omega|}$, and the distribution $\nu$ satisfies Slater's condition:
\[
\nu\in(0,\infty)^{|\Omega|},\qquad\mathbf1^\top\nu=1,\qquad a_e^\top\nu-\delta<0\qquad(e\in E(H)).
\]
We may therefore apply the ordinary finite-dimensional KKT multiplier
theorem, restricting the domain of $-\mathcal H_q$ to the open set $U:=(0,\infty)^{|\Omega|}$. Positivity is built into $U$, so the nonnegativity constraints need not be included. The KKT theorem yields numbers
\[
\lambda_e\geq0,\qquad\eta\in\mathbb R\qquad(e\in E(H))
\]
such that
\[
\lambda_e\bigl(a_e^\top\mu-\delta)=0\qquad (e\in E(H)),
\]
and
\[
\log_q(\mu_x)+\frac{1}{\ln q}+\sum_{e\in E(H)}\lambda_e(a_e)_x+\eta=0\qquad (x\in\Omega).
\]

\noindent Exponentiating,
\[
\mu_x=q^{-\sum_{e\in E(H)}\lambda_e(a_e)_x}q^{-\frac{1}{\ln q}-\eta}.
\]
Since 
\[
1=\sum_{x\in\Omega}\mu_x=q^{-\frac{1}{\ln q}-\eta}\left(\sum_{x\in\Omega}q^{-\sum_{e\in E(H)}\lambda_e(a_e)_x}\right),
\]
we can write 
\[
\mu_x=\frac{q^{-\sum_{e\in E(H)}\lambda_e(a_e)_x}}{Z(\lambda)},\ \text{where }Z(\lambda)=\sum_{x\in\Omega}q^{-\sum_{e\in E(H)}\lambda_e(a_e)_x}.
\]
Translating back to the logarithm we have
\[
-\log_q(\mu_x)=\sum_{e\in E(H)}\lambda_e(a_e)_x+\log_qZ(\lambda).
\]

We have found $\mu_x$ and $\log_q(\mu_x)$, so we may now find $\mathcal H_{q}(\mu)$:
\begin{align*}
    \mathcal H_{q}(\mu)&=-\sum_{x\in\Omega}\mu_x\log_q\mu_x\\
    &=\sum_{x\in\Omega}\mu_x\left(\sum_{e\in E(H)}\lambda_e(a_e)_x+\log_qZ(\lambda)\right)\\
    &=\left(\sum_{x\in\Omega}\mu_x\sum_{e\in E(H)}\lambda_e(a_e)_x\right)+\left(\sum_{x\in\Omega}\mu_x\cdot\log_q Z(\lambda)\right)\\
    &=\sum_{e\in E(H)}\lambda_ea_e^\top\mu+\log_qZ(\lambda)\\
    &=\delta\sum_{e\in E(H)}\lambda_e+\log_qZ(\lambda) \tag{1},
    \end{align*}
since by complementary slackness we have $\lambda_e(a_e^\top\mu-\delta)=0$ for each $e\in E(H)$.













Observe for any labeling $x$ and edge $e=uv$ we can write 
\[
(q^{-\lambda_e})^{\mathbf1_{x_u\not=x_v}}=q^{-\lambda_e}+(1-q^{-\lambda_e})\mathbf 1_{x_u=x_v},
\]

thus
\begin{align*}
    Z(\lambda)&=\sum_{x\in\Omega}\prod_{e\in E(H)}(q^{-\lambda_e}+(1-q^{-\lambda_e})\mathbf 1_{x_u=x_v})\\
    &=\sum_{x\in\Omega}\sum_{F\subseteq E(H)}\left((\prod_{e\in F}(1-q^{-\lambda_e})\mathbf 1_{x_u=x_v}))(\prod_{e\not\in F} q^{-\lambda_e})\right)\text{ by distribution laws.}\\
    &=\sum_{F\subseteq E(H)}\left((\prod_{e\in F}1-q^{-\lambda_e})(\prod_{e\not\in F} q^{-\lambda_e})\right)\sum_{x\in\Omega}\sum_{F\subseteq E(H)}\prod_{e\in F}\mathbf 1_{x_u=x_v}
\end{align*}

where
\[
\sum_{x\in\Omega}\sum_{F\subseteq E(H)}\prod_{e\in F}\mathbf 1_{x_u=x_v}=\sum_{F\subseteq E(H)}\sum_{x\in\Omega}\mathbf 1_{x_u=x_v\text{ for all }uv\in F}\geq\sum_{F\subseteq E(H)}q^{v(H)-|F|-1},
\]
since there are $q$ such choices for each of the $v(H)-1$ coordinates of $x$, except that coordinates in the same component of the graph defined by the edges of $F$ must be the same, and that graph surely has at most $|F|$ components.

We thus obtain
\begin{align*}
    Z(\lambda)&\geq\sum_{F\subseteq E(H)}\left(q^{v(H)-|F|-1}(\prod_{e\in F}1-q^{-\lambda_e})(\prod_{e\not\in F} q^{-\lambda_e})\right)\\
    &=q^{v(H)-1}\prod_{e\in E(H)}(q^{-\lambda_e}+\frac{1-q^{-\lambda_e}}{q})\text{ by distribution laws}.
\end{align*}

As a result we clearly have
\[
\log_q Z(\lambda)\geq v-1-m+\sum_e \log_q\bigl(1+(q-1)q^{-\lambda_e}\bigr),
\]

and so applying (1) we have shown
\[
\Phi_{H,q}(\delta)\geq v-1-m+\sum_e g_{q,\delta}(\lambda_e)\tag{2}
\]
where
\[
g_{q,\delta}(\lambda)=\log_q\bigl(1+(q-1)q^{-\lambda}\bigr)+\delta\lambda.
\]




We now minimize this one-variable function. Its derivative is
\[
g_{q,\delta}'(\lambda)=\delta-\frac{(q-1)q^{-\lambda}}{1+(q-1)q^{-\lambda}}.
\]

The unique critical point satisfies
\[
q^{-\lambda}=\frac{\delta}{(q-1)(1-\delta)}.
\]

Now since $\delta<1-\frac1q$, the right-hand side of the above expression is less than $1$, so this critical point lies in
$\lambda>0$. Furthermore,
\[
g_{q,\delta}''(\lambda)>0,
\]
so it is the minimum. At this point $\lambda_0$,
\[
1+(q-1)q^{-\lambda_0}=\frac1{1-\delta}
\]
and
\[
\lambda_0=\log_q\left(\frac{(q-1)(1-\delta)}{\delta}\right).
\]

\noindent Therefore, 
\begin{align*}
\min_{\lambda\geq0}g_{q,\delta}(\lambda)
&=-\log_q(1-\delta)+\delta\log_q\left(\frac{(q-1)(1-\delta)}{\delta}\right)\\
&=-(1-\delta)\log_q(1-\delta)-\delta\log_q\left(\frac{\delta}{q-1}\right)\\
&=H_q(\delta).
\end{align*}

\noindent Thus, every term in (2) is at least $H_q(\delta)$, proving
\[
\Phi_{H,q}(\delta)\geq v-1-m+mH_q(\delta).
\]

\noindent This completes the proof of the lemma.
\end{proof}

\end{document}
